% Standalone introduction section for inclusion in the JFM LaTeX template.
% The bibliography keys used here are defined in references.bib.

\section{Introduction}\label{sec:introduction}

The Boussinesq approximation accommodates small thermally induced density
variations within the incompressible Navier--Stokes equations by retaining them
only where they generate buoyancy.  When buoyancy is gravitational, the
underlying acceleration is prescribed and the usual reduction is comparatively
unambiguous \citep{GrayGiorgini1976}.  In a rapidly rotating flow, however, the
centrifugal acceleration can greatly exceed gravity, so even a small density
variation can reorganise the momentum balance.  This issue is especially
relevant to the rotating-disc boundary layer and the rotor--stator cavity,
whose basic states descend from the classical similarity solutions of
\citet{Cochran1934} and \citet{Batchelor1951} and underpin rotating-boundary-layer
stability theory \citep{Lingwood1995}.  Thermal extensions of these flows have
therefore commonly remained incompressible while introducing density variation
through prescribed centrifugal or rotating-force terms.  Such Boussinesq-type
closures have been used in similarity analyses of heated coaxial discs
\citep{Soong1996IJHMT,Soong1996IJRM}, in direct and large-eddy simulations of
rotor--stator heat transfer \citep{SerreEtAl2004,PoncetSerre2009,WangEtAl2023},
and in recent studies of thermally modified rotating-disc transition
\citep{DuEtAl2024}; see also the rotating-cavity review of
\citet{OwenLong2015}.

Unlike uniform gravity, centrifugal-type acceleration may arise from both
reference-frame rotation and the fluid's own swirl.  Retaining the density
perturbation only in a preselected solid-body centrifugal term is therefore a
modelling choice rather than a unique consequence of the Boussinesq limit.
\citet{LopezEtAl2013} showed that commonly used formulations omit centrifugal
buoyancy associated with differential rotation and strong vortices, while
\citet{BlackburnEtAl2021} derived a canonical formulation in which density
fluctuations interact consistently with the relevant accelerative terms.  The
latter analysis also showed that Coriolis and local-acceleration contributions
cannot always be discarded.  Centrifugal buoyancy is known to alter the onset
of rotating-cavity convection \citep{PitzEtAl2017}, and comparisons with fully
compressible calculations have revealed appreciable errors in reduced
rotating-cavity models \citep{GaoEtAl2020}.  Fully compressible rotating-disc
theory provides the corresponding reference once the small-density-variation
limit itself is no longer controlled \citep{TurkyilmazogluUygun2006}.
Nevertheless, model validity has usually been assessed through differences in
profiles, heat transfer or critical parameters.  Whether the closure can
change the multiplicity, topology and even existence of the basic state has
received much less attention.

Here we examine this more elementary form of model selectivity.  We first
continue the conventional centrifugal-Boussinesq equations from the isothermal
von K\'arm\'an state on the true semi-infinite domain and compare their behaviour
with an acceleration-consistent generalised Boussinesq formulation and the
fully compressible model.  We then continue the conventional closure in a
rotor--stator similarity system to determine whether finite-gap confinement,
the radial pressure-gradient constraint and the rotating core remove the
singularity found above a single disc.  Uniform disc temperatures are used, for
which the governing equations remain compatible with the similarity ansatz;
loss of a branch is therefore not caused by imposing a radially inconsistent
thermal boundary condition.  Pseudo-arclength continuation, a rational
Chebyshev representation of the semi-infinite domain and temporal spectra are
used to distinguish a dynamical fold from finite-domain or iterative-solver
failure.  The model hierarchy is interpreted asymptotically: the generalised
Boussinesq equations are tested as a small-temperature-difference approximation,
whereas the compressible equations are retained as the reference when density
linearisation is no longer justified.

The conventional model produces a non-degenerate saddle-node at
$T_{w,c}=1.048021731$ on the stable, isothermal-connected von K\'arm\'an branch.
A simple zero mode, stability exchange and square-root scaling confirm that
this premature termination is a property of the reduced equations, rather than
a loss of numerical convergence.  The acceleration-consistent basic state
continues through the same interval and its neutral curves agree closely with
the compressible results at small temperature differences, but the two models
diverge as density linearisation and constant-property assumptions lose their
asymptotic basis.  In that regime the fully compressible description is
necessary, rather than merely a more expensive refinement.  Rotor--stator
confinement delays, but does not generically remove, the topological failure:
at $Re_h=1000$ the conventional system has folds at $T_w=1.155517549$ and
$1.167676484$, with three similarity states between them.  These results do not
imply that every finite-radius, three-dimensional Boussinesq simulation must
fail at a universal temperature ratio, since radial non-similarity, sidewalls
and unsteadiness provide additional dynamics.  They do establish that the
choice of centrifugal-buoyancy closure is not merely a quantitative correction;
it can determine whether the basic state used for stability analysis exists at
all.
